% =============================================================================
% AGENT   : Lead Scientist + Geneticist (Domain Expert)
% MODULE  : Introduction — English version (intro_en.tex)
% SYNCED  : Must be consistent with intro_pt.tex (Lead Scientist responsibility)
% STATUS  : DRAFT
% =============================================================================

\section{Introduction}
\label{sec:intro}

The characterisation of genetic diversity is a cornerstone of plant
breeding, conservation genetics, and evolutionary biology
\autocite{Hamrick1989}. Molecular markers have become indispensable
tools for evaluating diversity at the genomic level, enabling
researchers to identify polymorphisms with high precision, reproducibility,
and throughput \autocite{Semagn2006}.

Among the marker technologies currently available, Single Nucleotide
Polymorphisms (\snp{}) and Simple Sequence Repeats (\ssr{}) are the
most widely employed. \snp{} markers represent the most abundant class
of genetic variation in any genome and are amenable to high-throughput
genotyping platforms such as SNP arrays and genotyping-by-sequencing
(GBS) \autocite{Elshire2011}. \ssr{} markers, also known as
microsatellites, are highly polymorphic, co-dominant, and have been
extensively validated in numerous crop species, making them valuable
for genetic mapping, parentage analysis, and variety identification
\autocite{Powell1996}.

The informativeness of a marker locus is commonly quantified using the
Polymorphism Information Content (\pic{}), first proposed by
\textcite{Botstein1980}. The \pic{} value ranges from 0 (monomorphic)
to a theoretical maximum approaching 0.5 for biallelic markers, with
values above 0.5 generally considered highly informative for multiallelic
loci. Expected (\he{}) and observed (\ho{}) heterozygosity are
complementary statistics that reflect the distribution of allelic
frequencies within a population and deviations from Hardy--Weinberg
equilibrium, respectively.

Despite extensive application of both marker classes in the literature,
comparative analyses of their discriminatory power under standardised
simulation conditions remain underexplored. Understanding the relative
performance of \snp{} and \ssr{} markers in terms of \pic{} and
heterozygosity is essential for designing efficient genotyping strategies
in future studies.

The present study aims to (i)~simulate a representative panel of
\snp{} and \ssr{} markers with biologically plausible allele frequency
distributions; (ii)~compute \pic{}, \he{}, and \ho{} for each marker
class; and (iii)~provide a visual and statistical comparison to inform
marker selection criteria in applied breeding programmes.
